% =====================================================================================================
% Section VII: Conclusion. Writer: conclusion writer, 2026-09-24. Numbers and sources: paper/notes/writer_conclusion.md
% =====================================================================================================

\section{Conclusion}\label{sec:concl}

Neural scaling laws are production functions, and fitting them is production-function estimation. Reading
them this way does three things. It organizes what scaling data can identify: cost-minimizing labs remove
the transverse variation that reveals curvature, so the better labs optimize, the less their data reveal
(Section~\ref{sec:ident}). It yields estimates with economic content: parameters and data are gross
complements in every public sweep, with an elasticity of substitution near 0.7 once the data are allowed to
choose the curvature (Section~\ref{sec:tech}). And it turns training choices into data: over-training reveals
anticipated inference demand, with a median open-weight base model built as if its lifetime inference compute
would be about 2.2 times its training compute (Section~\ref{sec:wedge}). The same lens disciplines
observational claims: a design-matched benchmark detects no bias in cross-lab returns to compute, while the
parameter--data split and the rate of ``algorithmic progress'' are weakly identified (Section~\ref{sec:obs}).
\textcolor{red}{[TBD-m9: one sentence on what our controlled two-lab experiment adds, e.g., $\sigma$ by lab with
$\kappa$ free and the neutrality verdict for FineWeb-Edu versus FineWeb.]}

\emph{What economists should use.}---For the substitutability of parameters and data, we recommend a range of
about 0.5 to 0.8 with a central value near 0.7: three estimators that do not impose the Chinchilla outer
exponent agree on 0.69--0.71 on the best-designed sweep, whereas the Chinchilla-form values of 0.73--0.83 rest
on a restriction that every sweep rejects. Complementarity itself is robust, because any interior compute
optimum requires $\sigma<1$ (Lemma~\ref{lemma:sigma}). For mapping compute into capability, the relevant object
is the frontier elasticity $\gamma$, which on-path data identify; its level varies across sweeps (about 0.10 to
0.18 under the Chinchilla form), so it should be taken from the sweep closest to the application. For the
deployment side, the wedge $w$ of Proposition~\ref{prop:wedge} measures how much of a model's lifetime compute
its developer planned to spend on inference. Raw exponents and the compute-optimal ratio $M^*$ are not portable:
$M^*$ at $10^{21}$ FLOP ranges from 3 to 60 across sweeps, and its 95 percent cluster-bootstrap interval at
Chinchilla's own compute is $[3.2,126]$.

\emph{Recommendations for scaling studies.}---Five practices follow from the identification results. First,
place a substantial share of runs off the expansion path and report the design statistic
$\mathrm{sd}(\ln M\mid\ln C)$: at equal total compute, an IsoFLOP design spanning 16 times the optimal model
size in each direction estimates $\sigma^*$ with a root-mean-squared error of 0.004, against 0.28 when every run
is compute-optimal (Figure~\ref{fig:designs}; Online Appendix~\ref{app:mc}). Second, test the outer exponent
$\kappa$ instead of imposing $\kappa=1$; the restriction is rejected in all seven sweeps, and relaxing it lowers
$\sigma^*$ by 0.04 to 0.28. Third, cluster inference by IsoFLOP budget: on the Chinchilla data, cluster-bootstrap standard errors
of the exponents are 1.5 to 2.2 times the pairs-bootstrap ones. Fourth, test the cross-equation restrictions
that link Approaches 1, 2, and 3 against a null that applies the same procedure to the fitted primal at the
observed design; naive comparisons reject even a correctly specified refit. Fifth, report the conventions behind
$N$ and $D$ and the uncertainty in $M^*$: counting conventions alone move the allocation exponent by 0.04 to
0.17 in the data we reanalyze. Two numerical points complete the list. Normalizing inputs at their geometric
means cuts the condition number of the Chinchilla Jacobian from about 2,000 to 62, and the Huber objective with
$\delta=10^{-3}$ on log loss is least absolute deviations in practice, with 84 percent of residuals in its
linear region, so inference should be derived for that estimator.

\emph{Implications for models of AI and growth.}---Integrated-assessment models of AI treat ``effective
compute'' as a sufficient statistic for capability \citep{erdil2025gate}, an aggregation that is exact only when
labs choose the input mix optimally, and \citet{trammell2023economic} note that scaling laws imply that compute and data are
gross complements. Our estimates quantify that complementarity and add two restrictions. Because $\sigma<1$, a
binding constraint on one input---for instance, a limited stock of text---lowers the return to the other faster
than a Cobb--Douglas aggregate would imply. And because the sign of biased
technical change is identified from the drift of compute-optimal tokens per parameter
(Proposition~\ref{prop:dmr}), published IsoFLOP laws over time can discipline models of directed technical
change \citep{acemoglu2002directed}. On the demand side, a model that counts only training compute understates
the planned lifetime compute of the median open-weight model by a factor of about three ($w=3.19$), with a range
of 1.9 to 4.0 across the technologies we consider.

\emph{Limitations.}---Four caveats bound these conclusions. First, units: losses come from different corpora and
tokenizers, so only unit-free objects compare across sweeps, and industry token counts are in each model's own
tokens. Second, extrapolation: three-quarters of the verified open-weight models exceed the largest compute
budget of any technology we use, so revealed demand is measured in ranks and ranges, not as precise levels
\textcolor{red}{[TBD-m9: direction of the extrapolation bias in $\hat w$ from our high-$M$ check]}. Third,
disclosure: developers that do not report training tokens are absent, and closed models rarely report them.
Fourth, functional form: the wedge and allocative calculations use the Chinchilla family or its generalization
with a free outer exponent, whereas a non-homothetic form fits the best sweep better out of sample. The
validation of revealed demand against usage is correlational.

\emph{Extensions.}---Two extensions follow directly. The first is the isoquant between training and test-time
compute. Repeated sampling and longer reasoning substitute inference computation for model quality
\citep{jones2021scaling,snell2024scaling,brown2024large}; the elasticity of substitution between the two is a
parameter that growth models of AI calibrate only loosely \citep{erdil2025gate}, and our duality and
identification arguments carry over to its estimation. The
second is market structure. The frontier implies a cost of capability that rises steeply as reducible loss
falls, $C^*\propto R^{-1/\gamma}$ with $1/\gamma\approx5.6$ under the Chinchilla refit; combined with demand that
escalates in quality, the endogenous-sunk-cost logic of \citet{sutton1991sunk} predicts a lower bound on
concentration that does not fall as the market grows \citep[see also][]{korinek2025concentrating}.

More broadly, the tools of empirical industrial organization apply directly to machine learning, and the
traffic runs both ways: designed training sweeps give production-function econometrics a rare setting in which
observational estimators can be benchmarked against experimental truth.
